\scanpage{32}
donc en vertu de $(1.1)$ on a
\[
 \xi_k=0\Longleftrightarrow\xi_{k'}=0
\]
où
\[
 k'=i+j'=\sum_0^r(p-1)\,p^\alpha=p^{r+1}-1,
\]
donc on trouve que
\[
 \xi_k=0\Longleftrightarrow\xi_{p^{r+1}-1}=0.
\]
Pour la $=$ raison, on trouve, si $\widetilde k\deleted{\unsure{=i+j'|=i+j}}\inserted{=i+\widetilde j}$ satisfait la $=$ condition que $k$, savoir $c_{i,\unsure{k'|\widetilde k}-i}=0$
\[
 \xi_{\widetilde k}=0\Longleftrightarrow\xi_{p^{r+1}-1}=0
\]
où on choisit $r$ tel que $p^{r+1}\geq\deleted{\unsure{k|i}}j,\widetilde j\ (\geq i)$.

En comparant, on trouve
\[
 \xi_k=0\Longleftrightarrow\xi_{\widetilde k}=0
\]
i.e. si \uline{un} $\xi_k$ avec $k\geq N$, $c_{k,k-i}\not\equiv0\ (p)$ est nul, tous le sont, cqfd.

\uline{Fin de la démonstration} Il \deleted{suffit de}\inserted{faut} prouver $\xi_k=0\quad\forall k\geq N$. On a
\[
 c_{l,k}\,\xi_k=c_{l,k-i}\,\xi_{k+l}\qquad\forall l\geq0
\]
et il suffit de prouver qu'il existe $l\geq0$ tel que $c_{l,k}\not\equiv0\ (p)$, et de plus $\xi_{k+l}=0$, ce qui est le cas \deleted{si} en vertu de $1.1.$ si $c_{i,(k+l)-i}\not\equiv0\ (p)$. Posons encore $k-i=j$ i.e. $k=i+j$, on veut
% archive sheet number: 31

